% Generated draft: requires mathematical and source review.
\documentclass[11pt]{article}
\usepackage[margin=25mm]{geometry}
\usepackage{fontspec}
\setmainfont{Latin Modern Roman}
\usepackage{amsmath,amssymb,mathtools,unicode-math}
\setmathfont{Latin Modern Math}
\usepackage{enumitem,needspace,xurl,hyperref}
\setlength{\parindent}{0pt}
\setlength{\parskip}{5pt}
\setlength{\emergencystretch}{3em}
\newcommand{\N}{\mathbb{N}}
\newcommand{\Z}{\mathbb{Z}}
\newcommand{\Q}{\mathbb{Q}}
\newcommand{\R}{\mathbb{R}}
\newcommand{\C}{\mathbb{C}}
\newcommand{\F}{\mathbb{F}}
\newcommand{\A}{\mathbb{A}}
\newcommand{\PP}{\mathbb{P}}
\newcommand{\E}{\mathbb{E}}
\newcommand{\eps}{\varepsilon}
\newcommand{\dd}{\,\mathrm{d}}
\begin{document}
\section*{UnsolvedMath ID: 30005480}
\textbf{Source problem code:} catalogue entry
\subsection{Definitions and mathematical statement}
For a nonzero polynomial \(f\in\R[t]\), real-rooted means that every complex root of \(f\) is real. Roots are counted with algebraic multiplicity whenever a root count is specified. If \(f(t)=a_m t^m+a_{m-1}t^{m-1}+\cdots\), with \(m\geq1\) and \(a_m\ne0\), then its roots have sum zero exactly when \(a_{m-1}=0\). A real polynomial \(g\) is a sum of squares (SOS) if there are finitely many real polynomials \(h_1,\ldots,h_r\) in the same variables such that \(g=\sum_{j=1}^r h_j^2\).

The supplied record requests the following two characterizations. Several essential terms are not defined in the supplied excerpt, so these are source-formulation schemata requiring clarification, rather than fully specified conjectures.

First, retain the source notation \(\R[t]_{n,0}\) and \(\R[t]_{d,0}\), and write
\[
q_n(t)=(t-1)^{n-1}(t+n-1).
\]
For every diagonal map
\[
T:\R[t]_{n,0}\longrightarrow\R[t]_{d,0},
\]
the requested assertion is
\[
\begin{gathered}
T(q_n)\text{ is real-rooted and has }d-1\text{ roots of one sign}
\\[2pt]
\Longleftrightarrow
\\[2pt]
T\text{ preserves zero-sum hyperbolicity}.
\end{gathered}
\]
Preservation here has the universally quantified form: every zero-sum hyperbolic input in the source domain has a zero-sum hyperbolic image in the source codomain. The excerpt does not define those domains, the bases with respect to which \(T\) is diagonal, or the intended meaning of zero-sum hyperbolicity. It also does not specify the admissible values of \(n,d\), degree or nonvanishing conditions on images, or whether the root condition means exactly or at least \(d-1\) roots, whether multiplicities count, and whether zero roots are allowed. The multiplicity convention above must therefore be checked against the intended source convention. The variable \(t\) replaces the source's \(x\) only as a dummy variable.

Second, for every hook-shaped symmetric polynomial \(p\) in the intended source class, the requested assertion is
\[
p\text{ is symmetrically hyperbolic}
\quad\Longleftrightarrow\quad
\Delta_{\mathbf 1,\mathbf 1}p\text{ is SOS}.
\]
Here \(\mathbf 1=(1,\ldots,1)\) in the ambient variable space. The supplied excerpt does not define hook-shaped, symmetrically hyperbolic, or \(\Delta_{u,v}p\), nor specify the variable count, degree, homogeneity, normalization, or other hypotheses on \(p\). These terms and hypotheses must be recovered before this equivalence can be treated as a precise mathematical conjecture.

There is also a material discrepancy between the requested characterizations and the supplied original excerpt. The latter states a quintic example as ``Theorem 9'' rather than stating either characterization. Let \(e_k\) denote the elementary symmetric polynomial of degree \(k\) in five variables. The excerpt's final monomial is transcribed as \(\texttt{e21 e3}\). If this means \(e_1^2e_3\), its displayed polynomial is
\[
p_{\mathrm{ex}}=4500e_5-220e_1e_4+7e_1^2e_3.
\]
This reading of the transcription requires confirmation. For the example, the excerpt asserts symmetric hyperbolicity, nonextendability of its associated operator, the SOS property of \(\Delta_{\mathbf 1,\mathbf 1}p_{\mathrm{ex}}\), and failure of SOS hyperbolicity. It attributes the SOS verification and a verification that \(\Delta_{\mathbf 1,u}p_{\mathrm{ex}}\) is not SOS for some \(u\in\R^5\) to computer calculations. It further records, for the associated operator \(T\),
\[
T\bigl((t-1)^4(t+4)\bigr)=750(t-1)^2(t-2)^2(t+6),
\]
and states that an associated operator satisfying
\[
T\bigl((t-1)^4(t+4)\bigr)=(t-a)^2(t-b)^2(t+2a+2b),
\qquad a,b\in\R_+,
\]
is not extendable. The meanings of associated operator, extendable, SOS hyperbolicity, and the boundary convention for \(\R_+\) are not supplied. These are attributed contextual assertions, not independently verified results or substitutes for the two requested characterizations.
\subsection{Short English statement}
The record requests two equivalences: a single test polynomial should characterize diagonal operators preserving zero-sum hyperbolicity, and an SOS condition should characterize symmetric hyperbolicity for hook-shaped symmetric polynomials. Essential definitions and parameter restrictions are missing. The supplied original excerpt instead discusses a quintic example, so the intended conjectures and their relationship to that excerpt require human mathematical review.
\subsection{Sources}
\begin{itemize}
\item Supplied local UnsolvedMath registry record, permanent ID \(30005480\), source problem code \texttt{catalogue entry}, titled \emph{Characterizations of Symmetrically Hyperbolic Hook-Shaped Polynomials}; collection: Oberwolfach Reports Open Problems. No individual proposer is identified in the supplied record.
\item Supplied report citation: \emph{New Directions in Real Algebraic Geometry} (2023), Oberwolfach Reports, pp.~813--881, DOI \texttt{10.4171/owr/2023/15}. \url{https://doi.org/10.4171/owr/2023/15}.
\item Supplied literature source: OWR 2023 report. \url{https://ems.press/content/serial-article-files/47009?nt=1}.
\item Supplied literature source: \emph{Symmetric Hyperbolic Polynomials}. \url{https://arxiv.org/abs/2308.09653}.
\item Supplied literature source: \emph{Journal of Pure and Applied Algebra}. \url{https://doi.org/10.1016/j.jpaa.2025.107869}.
\item Supplied dated literature curation, checked \texttt{2026-08-22}: the record labels the problem open and reports that the workshop report identifies an open problem, with evidence and low-degree or special-case results but no general characterization. This is supplied dated curation, not an independently verified assessment of the literature or current status.
\item The registry labels the statement \texttt{corrected\_verified} and says it was reconstructed and checked against the official report. This is a supplied verification claim; no source was independently visited or checked in preparing this draft, and the supplied original excerpt does not itself state the two requested equivalences.
\end{itemize}
\end{document}
